%% General definitions
\documentclass{article} %% Determines the general format.
\usepackage{a4wide} %% paper size: A4.
\usepackage[utf8]{inputenc} %% This file is written in UTF-8.
%% Some editors on Windows cannot save files in UTF-8.
%% If there is a problem with special characters not showing up
%% correctly, try switching "utf8" to "latin1" (ISO 8859-1).
\usepackage[T1]{fontenc} %% Format of the resulting PDF file.
\usepackage{fancyhdr} %% Package to create a header on each page.
\usepackage{lastpage} %% Used for "Page X of Y" in the header.
                      %% For this to work, you have to call pdflatex twice.
\usepackage{enumerate} %% Used to change the style of enumerations (see below).

\usepackage{amssymb} %% Definitions for math symbols.
\usepackage{amsmath} %% Definitions for math symbols.

\usepackage{tikz}  %% Pagacke to create graphics (graphs, automata, etc.)
\usetikzlibrary{shapes,arrows,calc,positioning}
\usetikzlibrary{graphs}
\usetikzlibrary{automata} %% Tikz library to draw automata
\usetikzlibrary{arrows}   %% Tikz library for nicer arrow heads
\usetikzlibrary{positioning}    %% Tikz library for more positioning options

\usepackage{hyperref} %% hyperlinks

\usepackage{chessfss} % contained in debian package texlive-games

\usepackage{listings}   % include & display python code
\usepackage[edges]{forest}  % for search trees

\tikzset{WhiteSquare/.style = {shape=rectangle, minimum width=28, minimum
    height=28, anchor=south west}}
\tikzset{BlackSquare/.style = {shape=rectangle, fill=lightgray, minimum
    width=28, minimum height=28, anchor=south west}}

%% Left side of header
\lhead{\course\\\semester\\Exercise \homeworkNumber}
%% Right side of header
\rhead{\authorname\\Page \thepage\ of \pageref{LastPage}}
%% Height of header
\usepackage[headheight=36pt]{geometry}
%% Page style that uses the header
\pagestyle{fancy}

\newcommand{\authorname}{Benjamin Regez\\Luis Tritschler}
\newcommand{\semester}{Spring Semester 2026}
\newcommand{\course}{Foundations of Artificial Intelligence}
\newcommand{\homeworkNumber}{9}

\renewcommand{\thesection}{Exercise \homeworkNumber.\arabic{section}}


\begin{document}

\section{Phase transitions for CNF formulas}
CNF formula $
    \varphi =
    \underbrace{(\ell_{11} \lor \ell_{12} \lor ... \lor \ell_{1k})}_{\text{clause } 1}
    \land
    \underbrace{(\ell_{21} \lor \ell_{22} \lor ... \lor \ell_{2k})}_{\text{clause } 2}
    \land
    ...
    \land
    \underbrace{(\ell_{m1} \lor \ell_{m2} \lor ... \lor \ell_{mk})}_{\text{clause } m}
$\\[1em]
with $n$ variables, $m$ clauses, and $k$ variables in each clause
\begin{enumerate}[(a)]
    \item Varying $k$ while fixing $n$ and $m$:\\[0.5em]
    For large values $k$, this will lead to an
    under-constrained region: For every
    additional variable in a clause, the clause
    has more "possibilities" to be satisfied and
    thus, with a large number of variables in each
    clause, the entire formula is more likely to be
    satisfied.
    \item Varying $n$ while fixing $m$ and $k$\\[0.5em]
    For large values $n$, this will lead to an
    over-constrained region: If the number of
    different variables occurring in the clauses
    increases, the number of variables that need
    to be set to "true" must also increase for
    the entire formula to be satisfied.
\end{enumerate}


\section*{Statement on scientific integrity}


\end{document}
